% GENERATED FILE. Edit canonical lessons, then run npm run generate:books.
% canonical-source-hash: fnv1a64:63266304cdc509d3
% canonical-lessons: GU-C315-udghatan, GU-C315-prasang, GU-C315-sthal, GU-C315-ullekh, GU-C315-nihalvu

\chapter{Opening Day}
\label{ch:gu-opening-day}
\hlchaptermodality{315}

\begin{chapteropening}
\textbf{By the end of this chapter:} \emph{I can say what happened on a special day and who watched.}

Complete the last lesson of chapter 315: I can say what happened on a special day and who watched.
\end{chapteropening}

\section[\texorpdfstring{\gu{ઉદ્ઘાટન} (udghāṭan)}{ઉદ્ઘાટન (udghāṭan)}]{\gu{ઉદ્ઘાટન} (udghāṭan) — an inauguration; a formal opening}

\label{lesson:GU-C315-udghatan}



\begin{quote}
Before the new one: say the Gujarati for a speech, then the Gujarati for to resound.
\end{quote}

\subsection*{You'll want to know: \gu{ઉદ્ઘાટન}}
\textbf{\gu{ઉદ્ઘાટન}} — \emph{udghāṭan} — \textquotedblleft{}an inauguration; a formal opening\textquotedblright{}.

\begin{quote}
\gu{આજે} \gu{નવા} \gu{પુલનું} \gu{ઉદ્ઘાટન} \gu{થયું}.
\end{quote}

\begin{quote}
\emph{āje navā pulnũ udghāṭan thayũ.}
\end{quote}

\begin{quote}
The new bridge was opened today.
\end{quote}

A neuter noun.

\begin{cousinweb}
\textbf{\gu{ઉદ્ઘાટન}} (\emph{udghāṭan}) is a tatsama, Sanskrit \emph{udghāṭana}, \textquotedblleft{}opening, unlocking\textquotedblright{} (Monier-Williams): the formal opening of a bridge, a road or a school.
\end{cousinweb}

\begin{grammarlens}[title={What you've built}]
One more word for telling what happened.
\end{grammarlens}

\subsection*{Your turn}
\begin{itemize}
\raggedright
  \item \emph{Say it:} \emph{udghāṭan}
  \item \emph{Say it:} \emph{udghāṭan}, once more
  \item \emph{Say it:} \emph{āje navā pulnũ udghāṭan thayũ.} — The new bridge was opened today.
  \item \emph{Recall:} say the Gujarati for a speech, then the Gujarati for to resound, then say \emph{udghāṭan} again
\end{itemize}

\begin{tcolorbox}[breakable,colback=teal!4,colframe=teal!35!black,title={Before you move on}]
What does \gu{ઉદ્ઘાટન} mean? (\textquotedblleft{}An inauguration; a formal opening\textquotedblright{}.) Say it once more.
\end{tcolorbox}
\section[\texorpdfstring{\gu{પ્રસંગ} (prasang)}{પ્રસંગ (prasang)}]{\gu{પ્રસંગ} (prasang) — an occasion}

\label{lesson:GU-C315-prasang}



\begin{quote}
Before the new one: say the Gujarati for to resound, then the Gujarati for an inauguration.
\end{quote}

\subsection*{You'll want to know: \gu{પ્રસંગ}}
\textbf{\gu{પ્રસંગ}} — \emph{prasang} — \textquotedblleft{}an occasion\textquotedblright{}.

\begin{quote}
\gu{આ} \gu{પ્રસંગે} \gu{ગામના} \gu{બધા} \gu{લોકો} \gu{ભેગા} \gu{થયા}.
\end{quote}

\begin{quote}
\emph{ā prasange gāmnā badhā loko bhegā thayā.}
\end{quote}

\begin{quote}
Everyone in the village gathered for the occasion.
\end{quote}

A masculine noun.

\begin{cousinweb}
\textbf{\gu{પ્રસંગ}} (\emph{prasang}) is a tatsama, Sanskrit \emph{prasaṅga}, \textquotedblleft{}connection\textquotedblright{}, and also \textquotedblleft{}an event, a case\textquotedblright{} (Monier-Williams). In Gujarati it is a special occasion: \textbf{\gu{આ} \gu{પ્રસંગે}}, on this occasion.
\end{cousinweb}

\begin{grammarlens}[title={What you've built}]
One more word for telling what happened.
\end{grammarlens}

\subsection*{Your turn}
\begin{itemize}
\raggedright
  \item \emph{Say it:} \emph{prasang}
  \item \emph{Say it:} \emph{prasang}, once more
  \item \emph{Say it:} \emph{ā prasange gāmnā badhā loko bhegā thayā.} — Everyone in the village gathered for the occasion.
  \item \emph{Recall:} say the Gujarati for to resound, then the Gujarati for an inauguration, then say \emph{prasang} again
\end{itemize}

\begin{tcolorbox}[breakable,colback=teal!4,colframe=teal!35!black,title={Before you move on}]
What does \gu{પ્રસંગ} mean? (\textquotedblleft{}An occasion\textquotedblright{}.) Say it once more.
\end{tcolorbox}
\section[\texorpdfstring{\gu{સ્થળ} (sthaḷ)}{સ્થળ (sthaḷ)}]{\gu{સ્થળ} (sthaḷ) — a place; a site}

\label{lesson:GU-C315-sthal}



\begin{quote}
Before the new one: say the Gujarati for an inauguration, then the Gujarati for an occasion.
\end{quote}

\subsection*{You'll want to know: \gu{સ્થળ}}
\textbf{\gu{સ્થળ}} — \emph{sthaḷ} — \textquotedblleft{}a place; a site\textquotedblright{}.

\begin{quote}
\gu{પત્રકારો} \gu{અકસ્માતના} \gu{સ્થળે} \gu{પહોંચ્યા}.
\end{quote}

\begin{quote}
\emph{patrakāro akasmātnā sthaḷe pahõcyā.}
\end{quote}

\begin{quote}
Journalists reached the scene of the accident.
\end{quote}

A neuter noun.

\begin{cousinweb}
\textbf{\gu{સ્થળ}} (\emph{sthaḷ}) is a tatsama, Sanskrit \emph{sthala}, \textquotedblleft{}ground, place, spot\textquotedblright{} (Monier-Williams). It is the formal word for the place where something happened: \textbf{\gu{અકસ્માતના} \gu{સ્થળે}}, at the scene.
\end{cousinweb}

\begin{grammarlens}[title={What you've built}]
One more word for telling what happened.
\end{grammarlens}

\subsection*{Your turn}
\begin{itemize}
\raggedright
  \item \emph{Say it:} \emph{sthaḷ}
  \item \emph{Say it:} \emph{sthaḷ}, once more
  \item \emph{Say it:} \emph{patrakāro akasmātnā sthaḷe pahõcyā.} — Journalists reached the scene of the accident.
  \item \emph{Recall:} say the Gujarati for an inauguration, then the Gujarati for an occasion, then say \emph{sthaḷ} again
\end{itemize}

\begin{tcolorbox}[breakable,colback=teal!4,colframe=teal!35!black,title={Before you move on}]
What does \gu{સ્થળ} mean? (\textquotedblleft{}A place; a site\textquotedblright{}.) Say it once more.
\end{tcolorbox}
\section[\texorpdfstring{\gu{ઉલ્લેખ} (ullekh)}{ઉલ્લેખ (ullekh)}]{\gu{ઉલ્લેખ} (ullekh) — a mention}

\label{lesson:GU-C315-ullekh}



\begin{quote}
Before the new one: say the Gujarati for an occasion, then the Gujarati for a site.
\end{quote}

\subsection*{You'll want to know: \gu{ઉલ્લેખ}}
\textbf{\gu{ઉલ્લેખ}} — \emph{ullekh} — \textquotedblleft{}a mention\textquotedblright{}.

\begin{quote}
\gu{છાપામાં} \gu{અમારા} \gu{ગામનો} \gu{ઉલ્લેખ} \gu{હતો}.
\end{quote}

\begin{quote}
\emph{chāpāmā̃ amārā gāmno ullekh hato.}
\end{quote}

\begin{quote}
Our village got a mention in the newspaper.
\end{quote}

A masculine noun.

\begin{cousinweb}
\textbf{\gu{ઉલ્લેખ}} (\emph{ullekh}) is a tatsama, Sanskrit \emph{ullekha}, \textquotedblleft{}mentioning, speaking of\textquotedblright{} (Monier-Williams). Something is mentioned with \textbf{\gu{કરવું}}: \textbf{\gu{ઉલ્લેખ} \gu{કર્યો}}.
\end{cousinweb}

\begin{grammarlens}[title={What you've built}]
One more word for telling what happened.
\end{grammarlens}

\subsection*{Your turn}
\begin{itemize}
\raggedright
  \item \emph{Say it:} \emph{ullekh}
  \item \emph{Say it:} \emph{ullekh}, once more
  \item \emph{Say it:} \emph{chāpāmā̃ amārā gāmno ullekh hato.} — Our village got a mention in the newspaper.
  \item \emph{Recall:} say the Gujarati for an occasion, then the Gujarati for a site, then say \emph{ullekh} again
\end{itemize}

\begin{tcolorbox}[breakable,colback=teal!4,colframe=teal!35!black,title={Before you move on}]
What does \gu{ઉલ્લેખ} mean? (\textquotedblleft{}A mention\textquotedblright{}.) Say it once more.
\end{tcolorbox}
\section[\texorpdfstring{\gu{નિહાળવું} (nihāḷvũ)}{નિહાળવું (nihāḷvũ)}]{\gu{નિહાળવું} (nihāḷvũ) — to watch; to behold}

\label{lesson:GU-C315-nihalvu}



\begin{quote}
Before the new one: say the Gujarati for a site, then the Gujarati for a mention.
\end{quote}

\subsection*{You'll want to know: \gu{નિહાળવું}}
\textbf{\gu{નિહાળવું}} — \emph{nihāḷvũ} — \textquotedblleft{}to watch; to behold\textquotedblright{}.

\begin{quote}
\gu{ઘણા} \gu{લોકોએ} \gu{ઉદ્ઘાટન} \gu{નિહાળ્યું}.
\end{quote}

\begin{quote}
\emph{ghaṇā lokoe udghāṭan nihāḷyũ.}
\end{quote}

\begin{quote}
Many people watched the opening.
\end{quote}

Past, with \gu{તેણે} (he or she …): \gu{તેણે} \gu{આકાશ} \gu{નિહાળ્યું}, \emph{teṇe ākāś nihāḷyũ}.

\begin{cousinweb}
\textbf{\gu{નિહાળવું}} (\emph{nihāḷvũ}) goes back to Sanskrit \emph{nibhālayati}, \textquotedblleft{}perceives, sees\textquotedblright{}, and CDIAL lists the Gujarati verb, \textquotedblleft{}to see\textquotedblright{} (CDIAL 7228). It is to watch with attention, a fuller word than \textbf{\gu{જોવું}}.
\end{cousinweb}

\begin{grammarlens}[title={What you've built}]
One more word for telling what happened.
\end{grammarlens}

\subsection*{Your turn}
\begin{itemize}
\raggedright
  \item \emph{Say it:} \emph{nihāḷvũ}
  \item \emph{Say it:} \emph{nihāḷvũ}, once more
  \item \emph{Say it:} \emph{ghaṇā lokoe udghāṭan nihāḷyũ.} — Many people watched the opening.
  \item \emph{Recall:} say the Gujarati for a site, then the Gujarati for a mention, then say \emph{nihāḷvũ} again
\end{itemize}

\begin{tcolorbox}[breakable,colback=teal!4,colframe=teal!35!black,title={Before you move on}]
What does \gu{નિહાળવું} mean? (\textquotedblleft{}To watch; to behold\textquotedblright{}.) Say it once more.
\end{tcolorbox}
